\section{Data preparation}
\label{sec:data-preparation}

\subsection{Thermodynamic conversion}

WBT is calculated from temperature $T$ and dew point $T_d$ in kelvin and
surface pressure $p$ in pascals. Vapour pressure is
$e(T_d)=611.2\exp\{17.67(T_d-273.15)/(T_d-29.65)\}$ Pa, and the Bolton
mixing-ratio, lifting-condensation-temperature and equivalent-potential-
temperature relations are evaluated in these units \citep{bolton1980}.
Saturated $T_w$ is obtained by solving
\begin{equation}
  \theta_e(T_w,T_w,p)=\theta_e(T,T_d,p)
  \label{eq:wbt}
\end{equation}
at the original pressure with 40 bisection iterations on
$[T_d,T]$; even a 100-K initial interval then falls below $10^{-9}$ K. Values
with $T_d>T$ are clipped at equality before evaluation. The Stull approximation
provides a comparison \citep{stull2011}; the main calculation follows the
pseudoadiabatic definition discussed by \citet{davies2008}.


\subsection{Quality control}

Every 16th longitude and 18th latitude cell from a fixed north-west origin
gives a $13\times13$ candidate lattice; removing ocean cells leaves 121 sites
(main-text Figure~2). The 35 summers contain 9,350,880
quality-controlled hour--site records and 92 complete UTC fields per summer.
Because of NetCDF packing, 2,092 dew-point values exceeded air temperature
slightly (0.022\% of records); these pairs were clipped at equality before WBT
was calculated.

\subsection{Alternative daily fields}

For each UTC day, the 121 simultaneous values at the peak regional mean form
the daily field. Within each month-year, the upper regional-mean quartile is
high, the middle half is middle and the lower quartile is excluded. Labels
depend on regional mean WBT, and the selected hour maximises that mean; graph
dispersion enters after selection. UTC+8 days,
alternative thresholds, sitewise maxima, all 24 fixed UTC hours and daily-mean
fields provide comparisons. At each fixed hour, both the original
peak-based labels and labels recomputed from that hour's regional mean are
analysed; daily-mean labels are likewise recomputed from daily-mean WBT\@. A
dated event manifest records the type-7 thresholds, membership and absence of
ties.

Three checks address within-month seasonal progression: comparison of the
day-of-month distributions, removal of a separate
linear day-of-month trend at every site within each year--month, and
subtraction of each site's calendar-day climatology, estimated from all
summers except the one being analysed. The last two calculations use the
original event labels.
A date-matched calculation compares each high day with
all middle days in the same month-year lying within $\pm3$ or $\pm5$ calendar
days. It first averages the available high-day-specific ratios and then uses
the same equal-month, equal-scale and equal-summer hierarchy as the
estimator; unmatched-day and empty-record counts are retained.

A continuous diagnostic regresses $\log Q_h$ on centred regional
mean WBT within each month-year and bandwidth. Slopes are then averaged equally
over months, bandwidths and evaluation summers. This slope supplements the
quartile target with a graded summary.


\subsection{Spatial weighting}

The main target is a site average. For area weighting,
$a_i=\cos(\mathrm{latitude}_i)$,
\[
 \bar y_t^{(a)}=\frac{\sum_i a_i y_{it}}{\sum_i a_i},\qquad
 Q_{h,a}(\mathbf y)=
 \frac{\sum_{i<j}a_i a_j w_{ij,h}(y_i-y_j)^2}
      {2\sum_{i<j}a_i a_j w_{ij,h}}.
\]
The first area-weighted calculation replaces $Q_h$ by $Q_{h,a}$ using the
original labels; the second redefines the labels from $\bar y_t^{(a)}$ at the
regional peak hours.
The two calculations separate spatial weighting from event relabelling.


Alternative graph kernels are $\exp(-d_{ij}/\rho)$ and
$[1-d_{ij}/\rho]_+^2$. For each Gaussian scale, $\rho$ is found by bisection so
that the alternative kernel has the same edge-weighted mean pair distance.
Across the five scales, their effective edge counts range from 334 to 7,112
and from 382 to 7,113, compared with 366 to 7,109 for the Gaussian graphs;
the full matched-distance and effective-edge metadata accompany the analysis.
